% !TEX program = xelatex
\documentclass[11pt,letterpaper]{article}
\usepackage[top=1in,bottom=1in,left=1.15in,right=1.15in]{geometry}
\usepackage{fontspec}
\usepackage{unicode-math}

\setmainfont{Latin Modern Roman}[
  Extension=.otf,
  UprightFont=lmroman10-regular,
  BoldFont=lmroman10-bold,
  ItalicFont=lmroman10-italic,
  BoldItalicFont=lmroman10-bolditalic,
]
\setsansfont{Latin Modern Sans}[
  Extension=.otf,
  UprightFont=lmsans10-regular,
  BoldFont=lmsans10-bold,
  ItalicFont=lmsans10-oblique,
  BoldItalicFont=lmsans10-boldoblique,
]
\setmonofont{Latin Modern Mono}[
  Extension=.otf,
  UprightFont=lmmono10-regular,
  ItalicFont=lmmono10-italic,
  Scale=0.85,
]

\usepackage{xcolor,titlesec,enumitem,fancyhdr,hyperref,xurl,ragged2e,microtype}
\usepackage{../ac-paper-essay}

\hypersetup{
  colorlinks=true,
  linkcolor=acpurple,
  urlcolor=acpurple,
  pdfauthor={@jeffrey},
  pdftitle={The Screen Looks Back},
}

\pagestyle{fancy}
\fancyhf{}
\renewcommand{\headrulewidth}{0pt}
\fancyhead[C]{\footnotesize\color{acgray}\thepage}

\newcommand{\ac}{Aesthetic\textcolor{acpink}{.}Computer}

\renewcommand{\essaytitle}[1]{%
  \begin{center}%
    \begin{tikzpicture}
      \node[acdark,opacity=0.25,
            font=\aclight\fontsize{38pt}{42pt}\selectfont,inner sep=0pt]
            at (2pt,-2pt) {#1};
      \node[acpink,
            font=\aclight\fontsize{38pt}{42pt}\selectfont,inner sep=0pt]
            at (0,0) {#1};
    \end{tikzpicture}%
  \end{center}%
  \vspace{0.6em}%
}

\begin{document}
\thispagestyle{empty}
\vspace*{0.45in}

\essaytitle{The Screen Looks Back}
\essaydate{August 5, 2026}
\begin{center}{\small\color{acgray}@jeffrey}\end{center}\par\vspace{1.4em}

A screen does two things at once. It shows you an image, and it records what you
do with it. For a long time the first action was obvious and the second could be
ignored. We called the screen an output device. The keyboard and mouse were
inputs. The computer sat politely between them, as if looking and being looked
at were separate jobs.

They are not separate anymore. The image changes because I touched it. The
system remembers where I touched. That memory determines the next image, which
invites another touch. At sufficient scale, the screen does not merely answer
me. It models me, predicts me, and prepares an answer before I have formed the
question.

Jan Tumlir and Jeffrey Stuker's \textit{From Medusa to Nvidia: On the
Algorithm's Primal Scene} gives this arrangement an ancient face. Their essay
begins with the linguistic crossing of algorithm, Algol, and the Gorgon, then
follows the image from Medusa's petrifying gaze through the mirror-shield,
photography, technical vision, search, artificial intelligence, and Nvidia's
spiraling eye. The route is deliberately allegorical. They do not explain the
algorithm as a procedure. They ask what kind of creature it has become in the
culture that uses it.\footnote{Jan Tumlir and Jeffrey Stuker, ``From Medusa to
Nvidia: On the Algorithm's Primal Scene,'' \textit{Artforum}, August 5, 2026:
\url{https://www.artforum.com/features/jan-tumlir-jeffrey-stuker-algorithm-1234756309/}.}

I make a computer whose entire public face is a screen. This is therefore not a
metaphor I can admire from a safe distance. It is a design review.

\section{The Shield Is Already a Screen}

Perseus survives Medusa by refusing a direct exchange of gazes. He looks into a
polished shield and acts through the reflection. In Tumlir and Stuker's account,
the shield is more than protection. It is an early image technology: a surface
that places representation between a body and a force too powerful to confront
without mediation.

The maneuver works, but it does not make the power disappear. Medusa's head is
cut off, carried away, and mounted on Athena's armor. Defense becomes weapon.
The image that saved one body becomes an instrument for arresting others. This
is the essay's most useful diagram for a computer screen. Mediation protects us
from the world by making it operable, but every operation can be reversed. The
screen through which I inspect the network is also the surface through which the
network inspects me.

Calling the screen a mirror understates the problem. A mirror does not ordinarily
keep a record. It does not compare my face with a billion others, infer what I
will do next, or auction the inference. The contemporary screen is mirror,
sensor, archive, and market joined into one plane. It reflects, but its
reflection has an owner.

This is why Nvidia is more than a topical flourish at the essay's end. The
company supplies machinery for rendering images, training models, simulating
worlds, and accelerating searches through the accumulated traces of other
images. Its eye-shaped mark gives a corporate body to the technical gaze. The
old power has not returned unchanged. It has become infrastructure.

\section{Aesthetic Computer Is Inside the Problem}

\ac{} opens to a prompt. Type a word and a piece fills the screen: a synthesizer,
a painting surface, a chat room, a game, a tiny Lisp program. The system's
immediate-mode renderer redraws the image every frame. A piece is not a static
page placed in front of a user; it exists in the exchange between code, screen,
speaker, and hand.\footnote{The present architecture and immediate-mode model
are documented in \textit{Aesthetic Computer '26}:
\url{https://papers.aesthetic.computer/aesthetic-computer-26-arxiv.pdf}.}

It would be easy to treat this interactivity as an automatic answer to the
Medusa problem. The person is doing something, not passively consuming. The
software is an instrument, not a feed. The code is open. There are no ads at the
front door and no recommender deciding which piece should arrive next.

Those differences matter. They do not constitute innocence.

The canvas still receives touch. Handles still connect actions to persistent
identities. Boot logs record device and performance facts. Piece logs collect
runtime events so that a failure on someone else's machine can be diagnosed.
Generated images can enter the same screen. A list can become a ranking; a
profile can become a performance for metrics; a useful log can become a habit of
collecting first and inventing reasons later.

Open source does not prevent capture. Neither does playfulness. A cheerful
instrument can gather more than it needs, and a transparent repository can sit
under an opaque social relation. Tumlir and Stuker's image is valuable here
because it moves the question away from declared intention. The issue is not
whether the maker means well. The issue is which direction the gaze travels and
what becomes possible after it lands.

\section{Address, Do Not Predict}

The strongest difference between \ac{} and an algorithmic feed is grammatical.
A feed predicts. A prompt waits.

On a feed, the system begins the sentence. It chooses an image on the basis of
earlier behavior, measures the response, and chooses again. The user may swipe,
like, hide, or leave, but every available gesture is already an answer to the
system's proposition. Even refusal becomes training data.

At the \ac{} prompt, the person begins with a name. Typing \texttt{notepat} or
visiting \texttt{aesthetic.computer/notepat} invokes that piece directly. The
address is stable, sayable, and independent of a personalized ranking. The
system does not need a model of the visitor to resolve the word.\footnote{This
design claim is developed in \textit{The URL Tradition}:
\url{https://papers.aesthetic.computer/url-tradition-26-arxiv.pdf}.}

This can sound like a minor interface preference. It is closer to a distribution
policy. Direct address prevents the system from quietly substituting prediction
for choice. It also limits discovery. A person can only type a name they know,
so lists, autocomplete, teaching, and conversation remain necessary. The honest
trade is between legibility and capture, not between a perfect prompt and a bad
feed.

The moment discovery becomes personalized, the Medusan loop returns. The system
must observe enough behavior to predict what should be shown; what it shows then
produces the behavior that validates the prediction. Preference hardens around
the instrument like stone. A creative system should expand the number of moves
a person can imagine, not become steadily more certain about the moves they will
make.

\section{Use, Do Not Hold}

An instrument wants an action. A feed wants duration.

This distinction is more reliable than the familiar division between active and
passive media. People act constantly inside feeds. They type, edit, perform,
argue, and make images. The platform converts all of that activity into a reason
to remain. Its primary measure is not the quality of the action but whether the
next action occurs on the same surface.

A musical instrument has another rhythm. I approach it, play, stop, and leave.
The silence after use is not a failure state. The instrument does not chase me
through the room with a better chord. A drawing tool can be finished with me. A
piece can end.

That is the standard \ac{} should hold itself to. A successful session need not
be long. A piece should not manufacture a reason to stay after its action is
complete. Metrics should diagnose the system rather than govern the work. If a
painting piece changes its brushes because a retention graph says red keeps
people three minutes longer, the instrument has become an attention machine even
if no advertisement appears.

The \textit{Sucking on the Complex} paper calls \ac{} an anti-environment: a
working structure built beside platform logic, capable of using a platform's
reach without inheriting its measurements.\footnote{\textit{Sucking on the
Complex}: \url{https://papers.aesthetic.computer/sucking-on-the-complex-26-arxiv.pdf}.}
Tumlir and Stuker sharpen that claim. An anti-environment is not defined by the
images it shows. It is defined by what those images are permitted to do to the
people looking at them.

\section{A Record Is a Design Decision}

Some observation is necessary. A networked instrument has to know enough to
connect players. A public creative system has to moderate abuse. A developer
needs evidence from a crash that happened on a phone across the world. The
answer cannot be to abolish memory and call the resulting blindness ethical.

The useful boundary is between a record that serves the action and a record that
turns the person into standing reserve.

A piece log has a named purpose: diagnose this run of this piece. Its fields can
be minimized, its retention bounded, and its audience restricted. Behavioral
prediction has no comparable stopping point. Every new signal might improve the
model, so the logic of collection expands until the person is represented by a
shadow larger than the action that cast it.

This means privacy cannot remain a policy written beside the interface. It has
to appear as limits in the system: events not emitted, content not copied,
identifiers not joined, histories allowed to expire. It should be possible to
say what each record is for without using the phrase ``improve the experience.''
The more specific the purpose, the more visible the moment when the record has
outlived it.

There is an aesthetic consequence. A system that refuses comprehensive memory
will sometimes seem less magical. It will fail to anticipate. It will ask again.
It will show the same front door to two different people. These are not merely
privacy costs. They are signs that the person has not been reduced to a
prediction.

\section{The Gaze in the Workshop}

Artificial image generation makes the problem literal. A model accepts a prompt
and returns a picture, but the apparent exchange conceals a much longer chain:
training images, annotators, model weights, data centers, accelerators, safety
filters, product interfaces, and the new prompt entering whatever records the
service retains. The image appears quickly because its history has been pushed
out of frame.

\ac{} can use these systems. So can its papers, posters, thumbnails, and pieces.
Refusal would not make the infrastructure cease to exist. Enthusiasm would not
make its provenance harmless. The practical task is to keep the generated image
from arriving as an oracle.

That means retaining its source conditions: which model made it, from which
instruction, for what use, with what human selection, and under what rights. It
also means refusing to confuse generation with completion. An image is inspected,
cropped, rejected, revised, connected to its consumer, or removed. The workshop
has to remain visible around the result.

Here the mirror-shield becomes useful again. Perseus does not defeat the gaze by
pretending he has no image technology. He uses a reflective surface deliberately,
at an angle, for a bounded action. The danger begins when the shield becomes a
permanent eye mounted on the institution.

\section{The Screen Looks Back}

Tumlir and Stuker end between acceleration and paralysis, with the phrase
``petrified unrest.'' The screen moves constantly while the body before it grows
still. Images mutate, refresh, and predict; the person supplies the attention
that keeps the movement going.

The opposite is not a motionless screen. It is a screen whose movement returns
to the person as capacity. After using it, can I do something I could not do
before? Can I play the pattern away from the device, remember the address, copy
the source, change the rule, invite someone else, or stop without penalty? If
the answer is yes, the image has behaved more like a score or an instrument than
a trap.

That is the relation between \textit{From Medusa to Nvidia} and \ac{}. The essay
does not supply a historical aura for the project. It supplies an adversarial
test. The prompt, the URL, the piece, the handle, the log, and the generated
image each have to answer the same question: does this surface increase the
person's agency, or use the appearance of agency to improve its picture of the
person?

A screen will look back. Sensors, networks, and shared computation guarantee as
much. The design question is whether its gaze is specific, temporary, and
answerable---whether I can see what it saw, understand why it mattered, and end
the exchange.

The screen that looks back must also let us look inside.

\colophon{\href{https://papers.aesthetic.computer/platter/readings/tumlir-stuker-medusa-nvidia-2026.txt}{Artforum source record and Aesthetic Computer research note}.}

\end{document}
